\chapter{NIST SP 800-22: Testes de Aleatoriedade para Geradores de Números Pseudoaleatórios}
Bateria de testes SP 800--22 do NIST e o conceito de indistinguibilidade estatística.

        O NIST Test Suite consiste em um pacote de 15 testes de estatistas desenvolvidos para avaliar a as sequencias binarias aleatórias (com comprimento arbitrário),
        gerados por geradores de números aleatórios ou pseudo criptografados, seja de forma em hardware ou software. 
        les podem ser decompostos em vários subtestes de acordo com a necessidade e o tipo de geração aleatória que podem ter uma sequencia. Os 15 testes são:
        \begin{enumerate}
          \item Teste de Frequência (monobit)
          \item Teste de Frequência dentro de um bloco
          \item Teste de Sequencias
          \item Testes para Maior Sequencia de Uns em um bloco
          \item Testes de Posto de Matriz Binária
          \item Teste da Transformada Discreta de Fourier (Espectral)
          \item Teste de Correspondência de Modelos não Sobrepostos
          \item Teste de Correspondência de Modelos Sobrepostos
          \item Teste "Estatístico Universal" de Maurer
          \item Teste de Complexidade Linear
          \item Teste Serial
          \item Teste de Entropia Aproximada
          \item Teste de Somas Cumulativas (Cusums)
          \item Teste de Excursões Aleatórias
          \item Teste de Variante de Excursões Aleatórias
        \end{enumerate}\nocite{NistsTP}

        Falar separadamente sobre cada um dos 15 testes. Mostrando um teste local separado. (Seguir documentacao do guthub para execucao local)\cite{NISTGit}
        
        \subsection{Bateria de Testes}%PASSO 1 - Explicacao de como vai ser efetuado os Testes

            Os arquivos binários gerados em Java servem como a entrada bruta para o NIST SP 800--22, 
            que interpreta os bytes como fluxos contínuos de bits para análise estatística. Ao executar o programa, 
            a ferramenta particiona esses dados em blocos de tamanho fixo, como sequências de $10^{6}$ bits, 
            aplicando sobre cada um deles quinze testes que buscam padrões não aleatórios, como repetições excessivas ou falta de complexidade linear. 
            Para cada teste, o sistema calcula um $P\text{-value}$, que quantifica a probabilidade de a sequência ter sido gerada por um processo puramente aleatório. 
            Ao final, o sucesso do gerador é determinado pela uniformidade desses valores e pela proporção de sequências que mantiveram o $P\text{-value} \geq 0,01$.

            \begin{figure}[htb]
                \centering
                \caption{Modelo simples de um PRNG com LCG}\label{fig:model_prng_lcg}
                \includegraphics[width=0.8\textwidth]{NIST/gerador_bits_testes_NIST.png}
                \legend{Fonte: Elaborado pelo autor (2026)}
            \end{figure}

        %PASSO 1.1 - Modelo de geracao LCG (voce já fez, tem pronto a explicacao)
        \subsubsection{Random}
        
        A função java.util.Random faz uso do modelo LCG “Gerador Congruencial Linear”. 
        Ele é um modelo puramente matemático. 
        Onde o mesmo é usado para gerar uma sequência de números aleatórios, 
        calculados por uma equação linear.
        \begin{equation}
            x_{i+1} = (A \cdot x_i + C) \pmod{M}
            \label{eq:lca}
        \end{equation}
        Onde o valor inicial “x0” é a seed, 
        A e C são parâmetros inteiros que podem ter um valor 
        já pré-definido ou passados diretamente, 
        M é o espaço dos valores aleatórios possíveis e “mod” é o 
        restante da divisão inteira.
        
        %PASSO 1.3 - Modelo de geracao Seguro usando entropia

        \subsubsection{SegureRandom}
        A função java.security.SegureRandom faz utilização do CSPRNG, que diferente do LCG.\@ Esse tipo de gerador busca caos do hardware.
        Sendo elas interrupçoes de hardware, pressionar das teclas do teclado, posição e movimento do mouse,
        pequenas flutuaçoes de voltagem do processador e/ou o tempo que o disco demora pra fazer a requisição.
        Esses eventos imprevisíveis são misturados, usando funçoes hash como SHA-256 para criar uma seed de alta entropia.
        Mesmo que se conheça o algoritmo, não é possivel recriar o mesmo evento. Justamente pelo fator caos fisico que foi utilizado.

        